% ABOUTME: Conclusion: what Phase I built and measured, where the project stands, and the Phase II plan.
% ABOUTME: Owner: Navodit. Every number here is from results/ (cpi.csv, performance.csv, synth.csv).
\section{Conclusion, Current Progress and Phase II Plan}
\label{sec:conclusion}

\subsection{Conclusion}
We built three processors for the same RV32I subset: a single-cycle CPU and
a five-stage pipeline with and without forwarding. Pipelining raised the
clock 1.94 times without forwarding and 1.43 times with it, but CPI rose to
1.51 to 2.74 and 1.26 to 1.58, so the net speedup ranged from 0.71 to 1.29.
The iron law explains each of these numbers: a pipeline is faster only when
the clock it gains is larger than the CPI it loses. In our programs the
largest single loss is the two-cycle flush after every taken branch, because
every loop iteration ends in one.

\subsection{Current progress}
Phase I is complete. All three CPUs are written in Verilog and pass all 248
automated checks against the reference simulator. They are synthesised and
placed for a Lattice ECP5 FPGA, and the CPI and speedup study is done. Every
table and figure in this report is generated from the repository by
\texttt{make synth perf figures}, so any change in Phase II can be measured
the same way.

\subsection{Phase II plan}
Phase II targets the two losses this report measures.
\begin{enumerate}
\item \emph{Branches.} Taken branch flushes add 0.23 to 0.35 to CPI. We will
first resolve branches in ID, which cuts the penalty to one cycle at the cost
of a comparator and forwarding paths in ID. Then we will add a dynamic
predictor with two-bit counters and a branch target buffer \cite{smith1981},
so a correctly predicted loop branch costs no cycles.
\item \emph{Stage balance.} EX is the longest stage in both pipelines (8.8 and
11.9\,ns). We will split it or move the branch decision out of it, so the
forwarding pipeline keeps more of its clock gain.
\item \emph{Instruction set.} Adding \texttt{jal}, \texttt{jalr} and byte and
halfword loads will let us run compiled C programs instead of hand-written
assembly.
\item \emph{Hardware.} We will run the design on a real ECP5 board using the
existing FPGA top level.
\end{enumerate}
Each change will be checked against the reference simulator and measured
with the same CPI and Fmax flow, so its gain can be compared directly with
the Phase I numbers.